\subsection{N-019: a viable maximal-scaling globalization and its complete table}
\label{sec:nonstandard-maximal-scaling}

\paragraph{A different, well-defined selection rule.}
There is a complete alternative to VIII's failed least-ambient-degree
selection.  Keep $K=\mathbb Q(v)$, $q=v^2$, $D=(v+v^{-1})^2$,
the actual orders $B_S$ over $S=\mathbb Q[D]$ and $B_R$ over
$R=\mathbb Q[v,v^{-1}]$, and the rational GT lines of N-010.
Fix the total GT label order recorded in the full Hermite certificate.
At each stage, with all lower exponents already fixed, choose the
\emph{largest} current exponent $a$ for which the external intersection
fiber is nonempty.  At that exponent the fiber is a singleton.
The resulting $114$ elements form an actual $S$-basis of $B_S$,
and an $R$-basis of $B_R$.  Their full polynomial multiplication
table is computed below.

This changes the selection criterion.  It is neither a claim that
VIII's ambient-degree comparator succeeds nor an identification with
its printed $E(4)$.  In fact the resulting basis fails the proposed
signed positivity property.  The construction supplies an explicit
viable alternative with a complete table and a precise positivity
boundary.  Unlike N-013's preliminary repair, its existence theorem
does not require imposing bar invariance on the candidate domain:
the output is automatically bar invariant.

\paragraph{Intrinsic Hermite description.}
Index the GT coordinates in \emph{descending} order by
$0,1,\ldots,n-1$, where $n=114$.  Thus the stage for coordinate $i$
has generic span of coordinates $i,i+1,\ldots,n-1$, and the induction
runs from $n-1$ down to zero.  Write $c'_i$ for its rational GT
matrix unit.  For each coordinate $j$ choose a monic common
denominator $\Delta_j(D)$ for the $j$-th coordinates of an actual
$S$-basis.  The resulting polynomial module has a unique row Hermite
basis $H_0,\ldots,H_{n-1}$, with GT coefficients
$h_{ij}\in F=\mathbb Q(D)$ satisfying
\[
 h_{ij}=0\quad(j<i),\qquad
 h_{ii}=p_i/\Delta_i,\quad p_i\text{ monic},\qquad
 \deg_D h_{ij}<\deg_D h_{jj}\quad(j>i).
 \tag{N-97}
\]
The last inequalities are the usual polynomial remainder conditions
after multiplying each column by its common denominator.  A larger
choice of denominator multiplies a column and its pivot by the same
polynomial; the quotient in polynomial division is unchanged.
The intrinsic selection proved below also shows denominator
independence directly.  Polynomial Euclidean row operations give
actual preimages, and reverse polynomial division proves equality
with the original module.  This is the full actual order, not the
saturation of its generic GT span.

Set
\[
 d_i=\deg_Dh_{ii},\qquad a_i=-d_i,\qquad
 c_i=q^{a_i}c'_i,\qquad
 L=\bigoplus_i A_0c_i,
 \tag{N-98}
\]
where $A_0\subset K$ consists of functions regular at $v=0$.
Since $D=q^{-1}(1+q)^2$, monicity implies
$h_{ii}/q^{a_i}=1\pmod {vA_0}$.
For $j>i$, the strict inequality in (N-97) implies
$h_{ij}/q^{a_j}\in vA_0$.
All $H_i$ belong to the actual $B_S$ and are bar fixed.
Consequently
\[
 H_i\in B_R\cap L\cap\bar L,
 \qquad H_i\longmapsto c_i\pmod {vL}.
 \tag{N-99}
\]

\paragraph{The triple is balanced.}
The projection is an isomorphism of rational vector spaces:
\[
 \begin{gathered}
 B_R\cap L\cap\bar L=\bigoplus_i\mathbb QH_i,\\
 B_R\cap L\cap\bar L\xrightarrow{\ \sim\ }L/vL.
 \end{gathered}
 \tag{N-100}
\]
To prove this without assuming bar invariance of an arbitrary element,
write it uniquely as $y=\sum_i f_i(v)H_i$, with $f_i\in R$;
actual base change and freeness are N-013/N-015.  Let $i$ be the
smallest index with $f_i\ne0$.  Its GT coordinate is $f_i h_{ii}$.
Membership in $L$ implies $\operatorname{ord}_{v=0}f_i\ge0$.
Membership in $\bar L$ implies $\deg_{v=\infty}f_i\le0$.
A Laurent polynomial satisfying both inequalities is a constant.
Subtract that rational constant times $H_i$, which belongs to both
lattices by (N-99), and continue.  This proves the displayed direct
sum.  Its projection sends the $H_i$ to the basis classes $c_i$,
so the projection is an isomorphism.  The same argument works in
every terminal GT flag subspace used in the induction.

\paragraph{Why the maximal current scaling selects exactly this basis.}
Fix stage $i$ and the lower exponents $a_j=-d_j$ for $j>i$.
Every actual element supported in this flag with nonzero current
coordinate has current GT coordinate $f(v)h_{ii}$ with nonzero $f\in R$.
Put $u=\operatorname{ord}_{v=0}f$ and $w=\deg_{v=\infty}f$,
so $u\le w$.  If it lies in the current fiber at exponent $a$,
its nonzero prescribed residue forces
\[
 u=2(a+d_i),\qquad
 w+2(a+d_i)\le0.
 \tag{N-101}
\]
The second inequality is the current-coordinate condition at infinity.
Since $w\ge u$, these imply $u\le0$ and therefore $a\le-d_i$.
At $a=-d_i$, one has $u=w=0$, so the residue normalization forces
$f=1$.  Equation (N-99) supplies the admissible element $H_i$.
Any two elements in this fiber differ by a lower-flag element in
$B_R\cap L\cap\bar L$ with zero residue, hence are equal by
(N-100).  Thus the largest exponent exists and its fiber is exactly
$\{H_i\}$.

This argument treats the entire external intersection fiber and all
permitted current exponents.  Over $K=\mathbb Q(v)$ a nonempty
fiber with a prescribed nonzero residue automatically has
$a\in\tfrac12\mathbb Z$; no missing rational exponent can produce
a larger value in this rational-order problem.  In the recorded
normalization every selected $a_i$ is integral.  Since the construction
agrees with the actual polynomial Hermite basis, it spans $B_S$ over
$S$ and not just its generic vector space.  Normalization of the same
GT lines by any nonzero $\mathbb Q(v)$ function transports the fibers
as in (N-96), translating their maximal exponents and multiplying
their unique elements by nonzero rational constants.  Thus it changes
neither the resulting actual-module basis property nor its existence.

\paragraph{The explicit rank-four data.}
The $114$ selected elements are exactly the already certified rows of
\nolinkurl{polynomial_order_gt_hermite.json.gz}.  Their GT coordinates,
all $300$ normalized opposite-KL coordinates, actual polynomial
preimages, and the complete total order are recorded there.
The independent full GT-map and actual-order verification is N-017.
The new table script additionally checks every remainder inequality
in (N-97).  It finds $H_0=1$ exactly.  The full vector of scaling
exponents is in \nolinkurl{gt_hermite_multiplication_verified.json}.
For the smallest $S'_{22}$ block they are $-8,-7,-6,-6$ in ascending
label order, distinguishing this selection from (N-84).

All $12996$ products are stored in
\nolinkurl{gt_hermite_multiplication_table.json.gz}.
The table has $1481544$ possible structure-constant slots, of which
$204876$ are nonzero; only $16$ ordered products vanish entirely.
Every constant belongs to $\mathbb Q[D]$ and its degree is at most
eight.  There are $130539$ distinct nonzero coefficient polynomials.
The exact construction is
\nolinkurl{gt_hermite_multiplication.py}.

Here is the complete algebraic verification mechanism.  Choose a
common polynomial denominator $\Delta$ for all rational GT
coordinates.  Write the matrix image of $H_i$ in the nine faithful
simple blocks as $N_i/\Delta$, with polynomial entries.
For every ordered pair, multiply these matrices exactly.  Actual
order closure implies, and the computation checks, that every
coordinate of $N_iN_j$ is divisible by $\Delta$.
Polynomial triangular division by the Hermite rows then gives
\[
 H_iH_j=\sum_k m_{ij}^k(D)H_k,
 \qquad m_{ij}^k(D)\in\mathbb Q[D].
 \tag{N-102}
\]
Every remainder in these divisions is checked to be zero.
The faithful GT map proves uniqueness and the equality in the
actual algebra.  Associativity is inherited from matrix
multiplication, not extrapolated from sample values or a bounded
set of triples.  The separate verifier
\nolinkurl{verify_gt_hermite_multiplication.py} imports no table
constructor and performs no Hermite divisions: it multiplies the
nine dense polynomial matrix blocks and checks the stored linear
combination in (N-102), after multiplying by $\Delta$.

\paragraph{A precise positivity obstruction remains.}
This viable basis is not a positive canonical basis in the sense
asked for in VIII.  With the certificate's zero-based row indices,
one exact structure constant is
\[
                 m_{1,10}^{92}(D)=D-\frac92.
 \tag{N-103}
\]
At $v=\sqrt2$ one has $D=9/2$, while $v-v^{-1}\ne0$.
Every nonzero Laurent polynomial with nonnegative coefficients is
strictly positive at that positive value of $v$.
Therefore (N-103) cannot have the form
$\pm(v-v^{-1})^a f(v)$ with $a\ge0$ and nonzero nonnegative
Laurent $f$, even without asking for unimodality.
This is an obstruction for this explicitly selected basis, not an
obstruction to every alternative basis or categorical construction.

For reproducibility, the table metadata also records the exact
signed Laurent tests: among the $130539$ distinct nonzero constants,
$52843$ have nonnegative coefficients after their forced sign and
maximal $(D-4)$ factor are removed, and $34203$ are also unimodal on
the integer $q$ lattice.  These counts describe the complete computed
table; the single coefficient (N-103) already proves the failure.
The full actual-order basis and full table remain valid regardless
of that positivity test.

\paragraph{Source and scope.}
The motivating external lattice/intersection problem is VIII,
section 3.3.  The maximal-scaling criterion, balanced-triple theorem,
and resulting explicit table are the separate construction here;
they are not attributed to VIII.  This is over the rational coefficient
order.  It does not contradict N-009's mixed-characteristic obstruction
to a $\mathbb Z[v,v^{-1}]$-basis.  No motivic realization, canonical
tensor category, or positive dimension ring is inferred from the
polynomial multiplication table alone.
